%% ===========================================================================
%%  APPENDICES.  Loaded after \appendices in main.tex.
%%
%%  Content that used to live here and has been PROMOTED to the main body:
%%    - Adaptive mask thresholding + Algorithm 1  -> Sec. IV-C
%%    - Complete editing pipeline, Algorithm 2    -> Sec. IV-D
%%    - Scale-level replacement analysis (Fig.)   -> Sec. VI-C
%% ===========================================================================

\section{Implementation Details}
\label{app:impl}

This appendix records the details needed to reproduce the pipeline exactly.

\noindent\textbf{Three generation passes, one set of weights.}
A single edit invokes the VAR transformer three times with identical weights and
seed. The passes differ only in four switches: whether continuous source features
are injected, which token storage is attached, whether the storage masks are
applied, and whether generated tokens are saved. Phase~1 injects the continuous
features at every scale with weight $0$, which replaces the accumulated residual
by the source features and makes the pass a reconstruction of $I_s$; Phase~2
attaches a storage holding the source bitwise tokens together with the preserve
mask, so background positions are pinned while the rest generates freely;
Phase~3 attaches the same tokens with the replacement mask.

\noindent\textbf{Which tokens are written back.}
The tokens substituted in Phase~3 are always the BSQ quantization of the source
image, obtained once before generation begins. The tokens produced by Phase~1 are
overwritten before Phase~3, and the tokens produced by Phase~2 are discarded.
Phase~2 therefore influences the output through exactly one channel, the target
focus mask --- which is the reason the preserve level $\ell$ is inert and the
reason removing Phase~2 altogether has a large effect
(Sec.~\ref{sec:ablation_phase2}).

\noindent\textbf{Mask polarity and granularity.}
Storage masks are \emph{replacement} masks, not protection masks: a true entry
means the position is copied from the source. Substitution operates on the
bitwise index $r_k \in \{0,1\}^{h_k \times w_k \times d}$ with $d{=}32$, so
composition is exact at the bit level and boundaries between replaced and
generated regions are resolved per bit.

\noindent\textbf{Attention extraction.}
Cross-attention is recorded by patching the forward pass of blocks $2$--$31$ and
recomputing $\mathrm{softmax}(qk^{\top}/\sqrt{d_h})$ from the module's own
weights, leaving the original output path untouched. Because VAR evaluates each
block exactly once per scale, the invocation counter of each block gives the
scale index directly. For a focus word spanning several text tokens, the
per-token maps are averaged; head dimensions are averaged before block-level
aggregation.

% TODO(optional): a data-flow figure of the three phases (mermaid source in
% docs/reports/p2p_three_phase_pipeline/report.md §0) would make this appendix
% considerably easier to follow.

\section{Evaluation Protocol}
\label{app:eval}

\noindent\textbf{Metrics.}
Structure Distance, and PSNR / SSIM / LPIPS / MSE restricted to the unedited
region, follow the official PIE-Bench implementation. CLIP whole and CLIP edited
use the standard CLIP score on the full image and on the annotated edited region
respectively. We additionally report HPSv2 (version 2.1) and ImageReward.

\noindent\textbf{Prompt convention.}
All preference and CLIP scores are computed against the \emph{original}
PIE-Bench editing prompt, with bracket annotations removed, regardless of which
prompt was used for generation. This matters: scoring a generation against a
paraphrased prompt instead of the original changes ImageReward by up to $0.33$
and HPSv2 by $0.025$ on this benchmark, which is larger than most of the effects
studied in Sec.~\ref{sec:analysis}. Any comparison that mixes the two conventions
is meaningless.

\noindent\textbf{Resolution handling.}
Edited images are resized to the source resolution with Lanczos resampling before
any metric is computed.

\noindent\textbf{Run-to-run noise.}
Generation uses bf16 kernels with fused attention and is not bit-exact. Repeating
a full-benchmark run under an identical configuration and seed changes the mean
by $\approx\!\pm 0.02$ ImageReward and $\approx\!\pm 0.001$ HPSv2 at $n{=}700$,
and by $\approx\!\pm 0.05$ ImageReward for an individual $n{=}80$ category. We
state this scale wherever a reported difference approaches it.

\section{Text-to-Image Editing Mode}
\label{app:t2i}

% TODO: move the T2I qualitative results here from the ACM supplementary
% material, and describe the mask-mechanism validation performed in this mode.

\section{Focus-Word Extraction: Failure Modes}
\label{app:prompt}

% TODO: write from docs/reports/m14_minimal_rewrite_pipeline (experiments 4-5)
% and docs/reports/prompt_length_comparison_m15.  Points to cover:
%   - 230/700 source prompts yield an empty focus set; concentrated in
%     add_object (75/80) and change_style (76/80).
%   - A conditional minimal patch (insert an anchor into the source prompt only,
%     leave the target prompt untouched) recovers +0.19 to +0.41 ImageReward on
%     the patched categories at no preservation cost.
%   - Rewriting whole prompts is strictly harmful and monotone in length
%     (IR 0.653 -> 0.372 -> 0.206 for original / short rewrite / long rewrite),
%     because generation follows the rewritten prompt while scoring uses the
%     original one.  This is the prompt-convention effect quantified in
%     Appendix~\ref{app:eval}, not a property of the editing method.
%   - Keep the main tables free of the rewriting module.

\section{Additional Qualitative Results}
\label{app:qualitative}

% TODO: 2-3 examples per editing category.  A journal appendix has the space and
% this is the cheapest way to answer "evaluation scope is narrow".
